\begin{definition}[\(\mathcal I\)-envelope]
\label{def:complex-i-envelope}
For \(\varphi\in\PSH(X,\theta)\), set
\[
 P_\theta[\varphi]_{\mathcal I}
 :=\uscsup\{\psi\in\PSH(X,\theta)\mid
          \psi\leq0,\ \psi\sim_{\mathcal I}\varphi\}.
\tag{51.27}\label{eq:complex-i-envelope}
\]
It is \(\mathcal I\)-model if it is fixed by this operator.
\end{definition}

\begin{theorem}\label{thm:complex-i-projection}
For every \(\varphi\in\PSH(X,\theta)\),
\[
 \varphi\sim_{\mathcal I}P_\theta[\varphi]_{\mathcal I},
\qquad
 P_\theta[P_\theta[\varphi]_{\mathcal I}]_{\mathcal I}
      =P_\theta[\varphi]_{\mathcal I}.
\tag{51.28}\label{eq:complex-i-projection}
\]
Furthermore,
\[
 P_\theta[\varphi]\leq P_\theta[\varphi]_{\mathcal I},
\qquad
 \varphi\sim_{\mathcal I}P_\theta[\varphi].
\tag{51.29}\label{eq:complex-p-i-comparison}
\]
If \(\varphi\) has analytic singularities, then
\[
 \varphi\sim P_\theta[\varphi]
        \sim P_\theta[\varphi]_{\mathcal I}.
\tag{51.30}\label{eq:complex-analytic-envelope-agreement}
\]
\end{theorem}

\begin{proof}
Choquet reduction replaces the family in
\eqref{eq:complex-i-envelope} by an increasing sequence.  Finite maxima
preserve divisorial Lelong numbers, and strong openness shows that its
usc supremum has the same multiplier ideals as \(\varphi\).  This proves
\eqref{eq:complex-i-projection}.  Formula
\eqref{eq:complex-p-envelope-rooftop}, again with strong openness, gives
\(\varphi\sim_{\mathcal I}P_\theta[\varphi]\), hence
\eqref{eq:complex-p-i-comparison}.  For analytic singularities, pass to
a log resolution.  Divisorial orders then control the whole singularity
up to \(O(1)\), so the two envelopes cannot be strictly less singular
than \(\varphi\).
\end{proof}
